\setcounter{chapter}{11}
\chapter{Dependency Declaration}\label{ch:12-dependency-declaration}

\chapterrule

In \hyperref[ch:06-local-dependencies]{Chapter 6}, we identified the layers the Notes API relies on: system libraries, the Python runtime, Python packages, and application configuration (plus its data). In \hyperref[ch:08-hidden-dependencies]{Chapter 8}, we saw how hidden dependencies~--- packages installed globally or tools assumed to be present~--- cause failures that are impossible to diagnose without the right mental model.

The first and most important step toward portability is making the Python package dependencies explicit. Explicit means written down in a file that travels with the source code. When someone (or something~--- a deployment system, a Docker build, a CI pipeline) needs to recreate the application's package environment from scratch, that file is the only instruction they need.

In Python, that file is \texttt{requirements.txt}.

\texttt{requirements.txt} is not glamorous. It is a plain text file with package names and version numbers. But it is the difference between a project that anyone can reproduce and a project that only works on the original developer's machine. It is the foundational document of Python project portability.

This chapter explains exactly what \texttt{requirements.txt} is, how to generate it, what version locking means and why it matters, and how pip uses it to recreate an environment. It also introduces \texttt{pyproject.toml}~--- the modern alternative~--- and explains when each is appropriate.

\begin{objectives}

By the end of this chapter, you will understand:

\begin{itemize}
\tightlist
\item
  What \texttt{requirements.txt} is and how it is structured
\item
  How to generate it from an active virtual environment
\item
  What version locking is, why it matters, and the difference between direct and transitive dependencies
\item
  How \texttt{pip\ install\ -}\allowbreak{}\texttt{r\ requirements.}\allowbreak{}\texttt{txt} recreates an environment
\item
  What \texttt{pip\ freeze} does and when to use it vs \texttt{pip\ list}
\item
  The difference between \texttt{requirements.txt} and \texttt{pyproject.toml} and when to use each
\end{itemize}

\end{objectives}

\setcounter{section}{0}
\section{\texorpdfstring{\texttt{requirements.txt}: The Manifest of What Your Program Needs}{requirements.txt: The Manifest of What Your Program Needs}}\label{12-dependency-declaration:121-requirementstxt-the-manifest-of-what-your-program-needs}

A \texttt{requirements.txt} file is a plain text list of Python packages and their version requirements. Here is a minimal one for the Notes API at this stage:

\begin{outputbox}[short]
\begin{Verbatim}[fontsize=\codesize, breaklines, breakanywhere, breaksymbolleft={\tiny\textcolor{muted}{$\hookrightarrow$}}]
# requirements.txt
flask==3.0.2
python-dotenv==1.0.1
\end{Verbatim}
\end{outputbox}

\noindent{}Each line specifies a package. The \texttt{==} operator specifies the exact version to install. This file travels with the source code in the repository. When someone sets up the project for the first time:

\begin{codebox}[short]

\begin{Shaded}
\begin{Highlighting}[]
\ExtensionTok{python3} \AttributeTok{{-}m}\NormalTok{ venv venv}
\BuiltInTok{source}\NormalTok{ venv/bin/activate}
\ExtensionTok{pip}\NormalTok{ install }\AttributeTok{{-}r}\NormalTok{ requirements.txt}
\end{Highlighting}
\end{Shaded}

\end{codebox}

\noindent{}They get these two packages at exactly these versions. They also get Flask's own dependencies~--- Werkzeug, Jinja2, and the rest~--- but those are \emph{not} pinned yet, so they may be different versions from yours. Section 12.2 explains why that matters, and how to pin them too.

\subsection*{\texorpdfstring{The structure of \texttt{requirements.txt}}{The structure of requirements.txt}}\phantomsection\label{12-dependency-declaration:the-structure-of-requirementstxt}

Each line in \texttt{requirements.txt} is a \textbf{requirement specifier}. The syntax supports several forms:

\begin{outputbox}
\begin{Verbatim}[fontsize=\codesize, breaklines, breakanywhere, breaksymbolleft={\tiny\textcolor{muted}{$\hookrightarrow$}}]
# Exact version (recommended for deployment)
flask==3.0.2

# Minimum version (allows newer versions)
flask>=3.0.0

# Compatible release (allows minor updates, ~= 3.0 means >= 3.0, < 4.0)
flask~=3.0

# Range constraint
flask>=3.0.0,<4.0.0

# No version constraint (get the latest — not recommended for production)
flask

# Comments (lines starting with #)
# This is a comment

# Include another file (typically used in requirements-dev.txt to pull in the base list)
-r requirements.txt

# Install from a local path
-e ./local-package

# Install from a git repository
git+https://github.com/org/repo.git@main
\end{Verbatim}
\end{outputbox}

\noindent{}For production deployments, exact versions (\texttt{==}) are strongly recommended. We explain why in Section 12.2.

\subsection*{\texorpdfstring{What goes in \texttt{requirements.txt}}{What goes in requirements.txt}}\phantomsection\label{12-dependency-declaration:what-goes-in-requirementstxt}

\texttt{requirements.txt} should list every package that \texttt{import} statements in your code rely on. These are your \textbf{direct dependencies}~--- the packages you explicitly use.

For the Notes API:

\begin{itemize}
\tightlist
\item
  \texttt{flask}: the web framework (direct dependency)
\item
  \texttt{python-dotenv}: for loading \texttt{.env} files (direct dependency)
\end{itemize}

\noindent{}What does not go in \texttt{requirements.txt}:

\begin{itemize}
\tightlist
\item
  \texttt{venv} or \texttt{pip} themselves (these are tools, not dependencies)
\item
  Packages that are part of Python's standard library (\texttt{sqlite3}, \texttt{os}, \texttt{json}, \texttt{sys}~--- these come with Python)
\item
  \texttt{\_\_pycache\_\_} or \texttt{.pyc} files (these are build artifacts, not dependencies)
\end{itemize}

\setcounter{section}{1}
\section{Version Locking: Why Exact Versions Matter}\label{12-dependency-declaration:122-version-locking-why-exact-versions-matter}

The \texttt{==} operator in \texttt{flask==3.0.2} specifies an exact version. This is called \textbf{version pinning} or \textbf{version locking}. It ensures that every time \texttt{pip\ install\ -}\allowbreak{}\texttt{r\ requirements.}\allowbreak{}\texttt{txt} runs, it installs exactly Flask 3.0.2~--- not 3.0.3, not 3.1.0, not 4.0.0. Always and only 3.0.2.

Why is this important?

\subsection*{Packages change between versions}\phantomsection\label{12-dependency-declaration:packages-change-between-versions}

Flask 3.0.2 and Flask 3.0.3 are very similar~--- 3.0.3 fixes bugs or makes small improvements. But even a patch release might:

\begin{itemize}
\tightlist
\item
  Change the default behavior of a feature you rely on
\item
  Deprecate an argument your code uses
\item
  Fix a bug that your code accidentally depended on
\end{itemize}

\noindent{}Without version pinning:

\begin{outputbox}[short]
\begin{Verbatim}[fontsize=\codesize, breaklines, breakanywhere, breaksymbolleft={\tiny\textcolor{muted}{$\hookrightarrow$}}]
# Without pinning, this might install Flask 3.0.2 today and a newer 3.x release tomorrow:
flask>=3.0.0
\end{Verbatim}
\end{outputbox}

\noindent{}When a colleague sets up the project a week later, they get a different version of Flask than you used. Their environment differs from yours. The ``but it worked on my machine'' problem returns, but now at the package version level.

\subsection*{Real-world consequences of unpinned versions}\phantomsection\label{12-dependency-declaration:real-world-consequences-of-unpinned-versions}

Consider this (hypothetical) scenario:

\begin{itemize}
\tightlist
\item
  Monday: Developer A installs Flask 3.0.2, builds and tests the API, everything works
\item
  Tuesday: a new patch release of Flask appears, fixing a security issue but also changing how \texttt{request.get\_json()} treats a malformed body
\item
  Wednesday: The CI/CD pipeline runs \texttt{pip\ install\ flask}, gets the new release, and tests the exact same code~--- but some tests now fail because \texttt{request.get\_json()} behavior changed
\end{itemize}

\noindent{}Now the CI is failing but the developer's machine passes. The versions are different. The investigation takes hours.

With pinned versions, Monday's \texttt{flask==3.0.2} is what everyone gets. The CI and the developer's machine behave identically.

\subsection*{Security updates and version pinning}\phantomsection\label{12-dependency-declaration:security-updates-and-version-pinning}

There is a tension here. Pinning versions prevents automatic updates, including security updates. If a newer release fixes a security vulnerability in 3.0.2, your pinned application continues to use the vulnerable version.

The resolution is: pin versions for stability, but periodically update your pins deliberately. This is a manual process (or automated with tools like \textbf{Dependabot} or \textbf{Renovate}): update \texttt{requirements.txt}, re-run tests, verify everything works, then deploy the update.

This is better than unpinned versions because the update is a conscious decision that goes through testing. Automatic unpinned updates can introduce breaking changes at any time.

\setcounter{section}{2}
\section{How pip Resolves and Installs Dependencies}\label{12-dependency-declaration:123-how-pip-resolves-and-installs-dependencies}

When you run \texttt{pip\ install\ flask}, pip does not just install Flask. Flask itself has dependencies, and those dependencies have dependencies. pip resolves the entire tree.

\subsection*{The dependency resolution process}\phantomsection\label{12-dependency-declaration:the-dependency-resolution-process}

\begin{codebox}[short]

\begin{Shaded}
\begin{Highlighting}[]
\ExtensionTok{pip}\NormalTok{ install flask}
\end{Highlighting}
\end{Shaded}

\end{codebox}

\noindent{}pip's internal resolver does the following:

\textbf{Step 1: Fetch Flask's metadata.} pip queries PyPI for Flask's metadata, which includes Flask's list of dependencies (specified in Flask's own \texttt{pyproject.toml}):

\begin{outputbox}[short]
\begin{Verbatim}[fontsize=\codesize, breaklines, breakanywhere, breaksymbolleft={\tiny\textcolor{muted}{$\hookrightarrow$}}]
Flask 3.0.2 requires:
  Werkzeug >= 3.0.0
  Jinja2 >= 3.1.2
  itsdangerous >= 2.1.2
  click >= 8.1.3
  blinker >= 1.6.2
\end{Verbatim}
\end{outputbox}

\noindent{}\textbf{Step 2: Fetch each dependency's metadata.} pip fetches the metadata for Werkzeug, Jinja2, etc. Each of those has its own dependencies (Jinja2 requires MarkupSafe, etc.).

\textbf{Step 3: Resolve conflicts.} If two packages require different versions of the same dependency, pip finds a version that satisfies both requirements. If no such version exists, pip reports a dependency conflict error.

\textbf{Step 4: Download wheel files.} pip downloads a compatible wheel for each package. Pure-Python packages (Flask, Jinja2, Werkzeug) have a single universal wheel (\texttt{py3-none-any}); packages with compiled code have one wheel per platform, and pip picks the one that matches your Python version, operating system, and CPU architecture.

\textbf{Step 5: Install.} pip extracts each wheel into the virtual environment's \texttt{site-packages}.

\subsection*{Direct vs transitive dependencies}\phantomsection\label{12-dependency-declaration:direct-vs-transitive-dependencies}

Your \texttt{requirements.txt} lists your \textbf{direct dependencies}~--- packages you explicitly import.

But when pip installs Flask, it also installs Flask's dependencies: Werkzeug, Jinja2, itsdangerous, click, blinker, and MarkupSafe. These are \textbf{transitive dependencies}~--- packages your code uses indirectly, through Flask.

Transitive dependencies are not listed in your \texttt{requirements.txt} (usually~--- see ``complete lockfiles'' below). They are discovered and installed automatically during the resolution process.

This creates a subtle problem: if you do not list transitive dependencies with their exact versions, they can change independently. Werkzeug 3.0.1 might be installed today and Werkzeug 3.0.3 a month from now~--- even though your \texttt{requirements.txt} only pins Flask.

\subsection*{\texorpdfstring{Generating a complete lockfile with \texttt{pip\ freeze}}{Generating a complete lockfile with pip freeze}}\phantomsection\label{12-dependency-declaration:generating-a-complete-lockfile-with-pip-freeze}

\textbf{\texttt{pip\ freeze}} prints the exact versions of every installed package~--- both direct and transitive:

\begin{codebox}[short]

\begin{Shaded}
\begin{Highlighting}[]
\ExtensionTok{pip}\NormalTok{ freeze}
\end{Highlighting}
\end{Shaded}

\end{codebox}

\noindent{}Output (after \texttt{pip\ install\ flask\ python-}\allowbreak{}\texttt{dotenv}):

\begin{outputbox}[short]
\begin{Verbatim}[fontsize=\codesize, breaklines, breakanywhere, breaksymbolleft={\tiny\textcolor{muted}{$\hookrightarrow$}}]
blinker==1.7.0
click==8.1.7
Flask==3.0.2
itsdangerous==2.1.2
Jinja2==3.1.3
MarkupSafe==2.1.5
python-dotenv==1.0.1
Werkzeug==3.0.1
\end{Verbatim}
\end{outputbox}

\noindent{}Saving this output to \texttt{requirements.txt} gives you a \textbf{complete lockfile}~--- every package at an exact version:

\begin{codebox}[short]

\begin{Shaded}
\begin{Highlighting}[]
\ExtensionTok{pip}\NormalTok{ freeze }\OperatorTok{\textgreater{}}\NormalTok{ requirements.txt}
\end{Highlighting}
\end{Shaded}

\end{codebox}

\noindent{}This is the most reproducible form of \texttt{requirements.txt}. Every package~--- direct and transitive~--- is pinned. Anyone who runs \texttt{pip\ install\ -}\allowbreak{}\texttt{r\ requirements.}\allowbreak{}\texttt{txt} gets an identical environment.

\subsection*{The two-file pattern}\phantomsection\label{12-dependency-declaration:the-two-file-pattern}

Many projects use two files:

\textbf{\texttt{requirements.txt}}: the complete lockfile from \texttt{pip\ freeze}, used for deployment. Every version is pinned.

\textbf{\texttt{requirements.in}} (or \texttt{requirements-base.txt}): only the direct dependencies, without pinning:

\begin{outputbox}[short]
\begin{Verbatim}[fontsize=\codesize, breaklines, breakanywhere, breaksymbolleft={\tiny\textcolor{muted}{$\hookrightarrow$}}]
# requirements.in
flask
python-dotenv
\end{Verbatim}
\end{outputbox}

\noindent{}Tools like \textbf{pip-tools} generate \texttt{requirements.txt} from \texttt{requirements.in}:

\begin{codebox}[short]

\begin{Shaded}
\begin{Highlighting}[]
\ExtensionTok{pip}\NormalTok{ install pip{-}tools}
\ExtensionTok{pip{-}compile}\NormalTok{ requirements.in  }\CommentTok{\# Generates requirements.txt with all transitive deps pinned}
\end{Highlighting}
\end{Shaded}

\end{codebox}

\noindent{}The advantage: you maintain the short list of direct dependencies (\texttt{requirements.in}), and pip-tools generates the complete reproducible lockfile (\texttt{requirements.txt}). Updating a package means changing \texttt{requirements.in} and regenerating \texttt{requirements.txt} with \texttt{pip-compile}.

\setcounter{section}{3}
\section{\texorpdfstring{Installing from \texttt{requirements.txt}}{Installing from requirements.txt}}\label{12-dependency-declaration:124-installing-from-requirementstxt}

To reproduce the environment from scratch on a new machine:

\begin{codebox}[short]

\begin{Shaded}
\begin{Highlighting}[]
\CommentTok{\# 1. Create and activate a virtual environment}
\ExtensionTok{python3} \AttributeTok{{-}m}\NormalTok{ venv venv}
\BuiltInTok{source}\NormalTok{ venv/bin/activate   }\CommentTok{\# macOS/Linux}

\CommentTok{\# 2. Install exactly what is listed}
\ExtensionTok{pip}\NormalTok{ install }\AttributeTok{{-}r}\NormalTok{ requirements.txt}
\end{Highlighting}
\end{Shaded}

\end{codebox}

\noindent{}\texttt{pip\ install\ -}\allowbreak{}\texttt{r\ requirements.}\allowbreak{}\texttt{txt} reads the file and installs every listed package at the specified version. It does nothing for packages already installed at the correct version~--- and it does not \emph{remove} packages that are installed but no longer listed (tools such as \texttt{pip-sync} from pip-tools do that).

\subsection*{\texorpdfstring{What \texttt{-r} means}{What -r means}}\phantomsection\label{12-dependency-declaration:what--r-means}

\texttt{-r} means ``requirements file''. It is the flag that tells pip to read a file rather than a single package name. You can combine multiple \texttt{-r} flags:

\begin{codebox}[short]

\begin{Shaded}
\begin{Highlighting}[]
\ExtensionTok{pip}\NormalTok{ install }\AttributeTok{{-}r}\NormalTok{ requirements.txt }\AttributeTok{{-}r}\NormalTok{ requirements{-}dev.txt}
\end{Highlighting}
\end{Shaded}

\end{codebox}

\subsection*{Verifying the installation}\phantomsection\label{12-dependency-declaration:verifying-the-installation}

After installing, verify the environment matches expectations:

\begin{codebox}[short]

\begin{Shaded}
\begin{Highlighting}[]
\ExtensionTok{pip}\NormalTok{ list}
\end{Highlighting}
\end{Shaded}

\end{codebox}

\noindent{}Or check specific packages:

\begin{codeboxcap}[short]{\textcolor{accent}{Listing 12.1}\enspace Verifying installed versions programmatically}

\begin{Shaded}
\begin{Highlighting}[]
\ImportTok{import}\NormalTok{ importlib.metadata}

\NormalTok{packages }\OperatorTok{=}\NormalTok{ [}\StringTok{"flask"}\NormalTok{, }\StringTok{"werkzeug"}\NormalTok{, }\StringTok{"python{-}dotenv"}\NormalTok{]}
\ControlFlowTok{for}\NormalTok{ pkg }\KeywordTok{in}\NormalTok{ packages:}
    \ControlFlowTok{try}\NormalTok{:}
\NormalTok{        version }\OperatorTok{=}\NormalTok{ importlib.metadata.version(pkg)}
        \BuiltInTok{print}\NormalTok{(}\SpecialStringTok{f"}\SpecialCharTok{\{}\NormalTok{pkg}\SpecialCharTok{\}}\SpecialStringTok{: }\SpecialCharTok{\{}\NormalTok{version}\SpecialCharTok{\}}\SpecialStringTok{"}\NormalTok{)}
    \ControlFlowTok{except}\NormalTok{ importlib.metadata.PackageNotFoundError:}
        \BuiltInTok{print}\NormalTok{(}\SpecialStringTok{f"}\SpecialCharTok{\{}\NormalTok{pkg}\SpecialCharTok{\}}\SpecialStringTok{: NOT INSTALLED"}\NormalTok{)}
\end{Highlighting}
\end{Shaded}

\end{codeboxcap}

\noindent{}Output:

\begin{outputbox}[short]
\begin{Verbatim}[fontsize=\codesize, breaklines, breakanywhere, breaksymbolleft={\tiny\textcolor{muted}{$\hookrightarrow$}}]
flask: 3.0.2
werkzeug: 3.0.1
python-dotenv: 1.0.1
\end{Verbatim}
\end{outputbox}

\setcounter{section}{4}
\section{\texorpdfstring{The Complete \texttt{requirements.txt} for the Notes API}{The Complete requirements.txt for the Notes API}}\label{12-dependency-declaration:125-the-complete-requirementstxt-for-the-notes-api}

At this stage of the book, the Notes API uses:

\begin{codebox}[short]

\begin{Shaded}
\begin{Highlighting}[]
\CommentTok{\# Direct imports in app.py:}
\ImportTok{from}\NormalTok{ flask }\ImportTok{import}\NormalTok{ Flask, request, jsonify, g   }\CommentTok{\# → flask}
\ImportTok{from}\NormalTok{ dotenv }\ImportTok{import}\NormalTok{ load\_dotenv                  }\CommentTok{\# → python{-}dotenv}
\ImportTok{import}\NormalTok{ sqlite3                                  }\CommentTok{\# standard library — no pip package}
\ImportTok{import}\NormalTok{ os                                       }\CommentTok{\# standard library}
\ImportTok{import}\NormalTok{ time                                     }\CommentTok{\# standard library}
\end{Highlighting}
\end{Shaded}

\end{codebox}

\noindent{}The direct dependencies are Flask and python-dotenv. After installing them in a fresh virtual environment and running \texttt{pip\ freeze}:

\begin{outputbox}[short]
\begin{Verbatim}[fontsize=\codesize, breaklines, breakanywhere, breaksymbolleft={\tiny\textcolor{muted}{$\hookrightarrow$}}]
# requirements.txt (generated by pip freeze after pip install flask python-dotenv)
blinker==1.7.0
click==8.1.7
Flask==3.0.2
itsdangerous==2.1.2
Jinja2==3.1.3
MarkupSafe==2.1.5
python-dotenv==1.0.1
Werkzeug==3.0.1
\end{Verbatim}
\end{outputbox}

\noindent{}Save this file at the root of the project, alongside \texttt{app.py}.

The complete project structure is now:

\begin{diagrambox}[short]{266.4pt}
\begin{Verbatim}[fontsize=\fontsize{9.0}{8.55}\selectfont]
notes-api/
├── app.py
├── requirements.txt
├── .env                   ← NOT committed
├── .env.example           ← committed (no real values)
├── .gitignore
└── venv/                  ← NOT committed
\end{Verbatim}
\end{diagrambox}

\noindent{}\texttt{.gitignore} should contain:

\begin{outputbox}[short]
\begin{Verbatim}[fontsize=\codesize, breaklines, breakanywhere, breaksymbolleft={\tiny\textcolor{muted}{$\hookrightarrow$}}]
# .gitignore
venv/
__pycache__/
*.pyc
*.db
.env
\end{Verbatim}
\end{outputbox}

\noindent{}With \texttt{requirements.txt} committed and \texttt{venv/} excluded, a collaborator can:

\begin{codebox}[short]

\begin{Shaded}
\begin{Highlighting}[]
\FunctionTok{git}\NormalTok{ clone https://github.com/yourname/notes{-}api.git}
\BuiltInTok{cd}\NormalTok{ notes{-}api}
\ExtensionTok{python3} \AttributeTok{{-}m}\NormalTok{ venv venv}
\BuiltInTok{source}\NormalTok{ venv/bin/activate}
\ExtensionTok{pip}\NormalTok{ install }\AttributeTok{{-}r}\NormalTok{ requirements.txt}
\FunctionTok{cp}\NormalTok{ .env.example .env}
\CommentTok{\# Edit .env with local values}
\ExtensionTok{python}\NormalTok{ app.py}
\end{Highlighting}
\end{Shaded}

\end{codebox}

\noindent{}This workflow is reproducible. No hidden global dependencies. No assumptions about what the collaborator has installed.

\setcounter{section}{5}
\section{\texorpdfstring{\texttt{pyproject.toml}: The Modern Alternative}{pyproject.toml: The Modern Alternative}}\label{12-dependency-declaration:126-pyprojecttoml-the-modern-alternative}

Python has a modern package configuration standard defined in \textbf{PEP 517}, \textbf{PEP 518}, and \textbf{PEP 621}, centered on a file called \texttt{pyproject.toml}. \texttt{pyproject.toml} is now the recommended way to declare Python project metadata, including dependencies.

For a simple application (as opposed to a library meant to be distributed), \texttt{requirements.txt} remains the most common and straightforward choice. But understanding \texttt{pyproject.toml} is valuable because it is increasingly seen in modern Python projects.

\begin{codebox}

\begin{Shaded}
\begin{Highlighting}[]
\CommentTok{\# pyproject.toml}
\KeywordTok{[project]}
\DataTypeTok{name} \OperatorTok{=} \StringTok{"notes{-}api"}
\DataTypeTok{version} \OperatorTok{=} \StringTok{"0.1.0"}
\DataTypeTok{description} \OperatorTok{=} \StringTok{"A simple notes API"}
\DataTypeTok{requires{-}python} \OperatorTok{=} \StringTok{"\textgreater{}=3.10"}

\DataTypeTok{dependencies} \OperatorTok{=} \OperatorTok{[}
    \StringTok{"flask\textgreater{}=3.0.0"}\OperatorTok{,}
    \StringTok{"python{-}dotenv\textgreater{}=1.0.0"}\OperatorTok{,}
\OperatorTok{]}

\KeywordTok{[project.optional{-}dependencies]}
\DataTypeTok{dev} \OperatorTok{=} \OperatorTok{[}
    \StringTok{"pytest"}\OperatorTok{,}
    \StringTok{"black"}\OperatorTok{,}
    \StringTok{"ruff"}\OperatorTok{,}
\OperatorTok{]}

\KeywordTok{[build{-}system]}
\DataTypeTok{requires} \OperatorTok{=} \OperatorTok{[}\StringTok{"hatchling"}\OperatorTok{]}
\DataTypeTok{build{-}backend} \OperatorTok{=} \StringTok{"hatchling.build"}

\KeywordTok{[tool.hatch.build.targets.wheel]}
\DataTypeTok{only{-}include} \OperatorTok{=} \OperatorTok{[}\StringTok{"app.py"}\OperatorTok{]}   \CommentTok{\# the project is still a single module (Chapter 15 makes it a package)}
\end{Highlighting}
\end{Shaded}

\end{codebox}

\noindent{}With \texttt{pyproject.toml}, you install the project with:

\begin{codebox}[short]

\begin{Shaded}
\begin{Highlighting}[]
\ExtensionTok{pip}\NormalTok{ install }\AttributeTok{{-}e}\NormalTok{ .          }\CommentTok{\# Installs the project and its dependencies}
\ExtensionTok{pip}\NormalTok{ install }\AttributeTok{{-}e} \StringTok{".[dev]"}   \CommentTok{\# Includes development dependencies too}
\end{Highlighting}
\end{Shaded}

\end{codebox}

\noindent{}The \texttt{-e} flag installs in \textbf{editable mode}~--- the source code is used directly, so changes to \texttt{app.py} take effect without reinstalling.

\Needspace{16\baselineskip}

\subsection*{\texorpdfstring{When to use \texttt{pyproject.toml} vs \texttt{requirements.txt}}{When to use pyproject.toml vs requirements.txt}}\phantomsection\label{12-dependency-declaration:when-to-use-pyprojecttoml-vs-requirementstxt}

{\def\LTcaptype{none} % do not increment counter
\begin{longtable}[]{@{}
  >{\raggedright\arraybackslash}p{(\linewidth - 2\tabcolsep) * \real{0.5000}}
  >{\raggedright\arraybackslash}p{(\linewidth - 2\tabcolsep) * \real{0.5000}}@{}}
\toprule\noalign{}
\begin{minipage}[b]{\linewidth}\raggedright
Use case
\end{minipage} & \begin{minipage}[b]{\linewidth}\raggedright
Recommendation
\end{minipage} \\
\midrule\noalign{}
\endhead
\bottomrule\noalign{}
\endlastfoot
Simple application (not a library) & \texttt{requirements.txt} with \texttt{pip\ freeze} \\
Application with dev/test/prod dependency groups & \texttt{pyproject.toml} + a lockfile tool (pip-tools, Poetry, uv) \\
A library that others will install & \texttt{pyproject.toml} is required \\
Modern project following current Python standards & \texttt{pyproject.toml} + lockfile \\
\end{longtable}
}

For the Notes API, we continue with \texttt{requirements.txt} for clarity. The concepts (exact version pinning, separating direct from transitive dependencies) are identical regardless of the file format.

\unnumberedsection{false}{The Updated Development Workflow}\label{12-dependency-declaration:the-updated-development-workflow}

With \texttt{requirements.txt} in place, the correct development workflow is:

\begin{codebox}

\begin{Shaded}
\begin{Highlighting}[]
\CommentTok{\# Setting up (first time on a machine)}
\ExtensionTok{python3} \AttributeTok{{-}m}\NormalTok{ venv venv}
\BuiltInTok{source}\NormalTok{ venv/bin/activate}
\ExtensionTok{pip}\NormalTok{ install }\AttributeTok{{-}r}\NormalTok{ requirements.txt}

\CommentTok{\# Adding a new package}
\ExtensionTok{pip}\NormalTok{ install gunicorn}
\ExtensionTok{pip}\NormalTok{ freeze }\OperatorTok{\textgreater{}}\NormalTok{ requirements.txt  }\CommentTok{\# Update the lockfile}
\FunctionTok{git}\NormalTok{ add requirements.txt}
\FunctionTok{git}\NormalTok{ commit }\AttributeTok{{-}m} \StringTok{"Add gunicorn to requirements"}

\CommentTok{\# After pulling changes that updated requirements.txt}
\ExtensionTok{pip}\NormalTok{ install }\AttributeTok{{-}r}\NormalTok{ requirements.txt  }\CommentTok{\# Update the environment to match}

\CommentTok{\# Updating a package deliberately}
\ExtensionTok{pip}\NormalTok{ install }\AttributeTok{{-}{-}upgrade}\NormalTok{ flask}
\ExtensionTok{pip}\NormalTok{ freeze }\OperatorTok{\textgreater{}}\NormalTok{ requirements.txt}
\CommentTok{\# Run tests to verify nothing broke}
\FunctionTok{git}\NormalTok{ add requirements.txt}
\FunctionTok{git}\NormalTok{ commit }\AttributeTok{{-}m} \StringTok{"Update Flask to 3.0.3"}
\end{Highlighting}
\end{Shaded}

\end{codebox}

\noindent{}The key habit: any time \texttt{requirements.txt} changes (because you added, removed, or updated a package), regenerate it with \texttt{pip\ freeze} and commit the result. The file should always reflect the actual installed state of the project's virtual environment.

This works only while the virtual environment contains \emph{nothing but} the application's runtime packages. As soon as you install development tools into it~--- a test runner, a linter (→ \hyperref[ch:17-build-process]{Chapter 17})~--- \texttt{pip\ freeze} would record them as production dependencies too. From that point on, keep development tools in a separate \texttt{requirements-dev.txt}, and maintain \texttt{requirements.txt} deliberately: edit it by hand, or generate it with a lock tool such as \texttt{pip-tools} (\texttt{pip-compile}) or \texttt{uv}.

\phantomsection\label{12-dependency-declaration:key-takeaways}
\begin{takeaways}

\begin{itemize}
\tightlist
\item
  \textbf{\texttt{requirements.txt}} is a text file that lists Python package dependencies and their versions. It makes the dependency environment reproducible by anyone, on any machine.
\item
  \textbf{Version pinning} (\texttt{flask==3.0.2}) ensures that \texttt{pip\ install\ -}\allowbreak{}\texttt{r\ requirements.}\allowbreak{}\texttt{txt} installs the same version every time, not whatever the latest version happens to be. Use \texttt{==} for production requirements.
\item
  Python packages have their own dependencies (\textbf{transitive dependencies}). \texttt{pip\ freeze} captures both direct and transitive dependencies at their exact versions. This is the most reproducible form of \texttt{requirements.txt}.
\item
  \texttt{pip\ install\ -}\allowbreak{}\texttt{r\ requirements.}\allowbreak{}\texttt{txt} installs every listed package at the specified version into the active virtual environment.
\item
  The \textbf{two-file pattern} (\texttt{requirements.in} for direct deps + \texttt{requirements.txt} for the full lockfile, generated by pip-tools) is a clean way to separate ``what I directly need'' from ``the full reproducible environment''.
\item
  \texttt{pyproject.toml} is the modern standard for project metadata and dependency declaration, especially for libraries. \texttt{requirements.txt} remains common for applications.
\item
  Add \texttt{requirements.txt} to version control. Never add \texttt{venv/}, \texttt{\_\_pycache\_\_/}, \texttt{.env}, or \texttt{.db} files.
\end{itemize}

\end{takeaways}

\unnumberedsection{false}{Exercises}\label{12-dependency-declaration:exercises}

\begin{enumerate}
\def\labelenumi{\arabic{enumi}.}
\tightlist
\item
  \textbf{Predict.} Compare \texttt{pip\ freeze} in your project environment with \texttt{requirements.txt}. Which packages appear in one and not the other, and why?
\item
  \textbf{Break and fix.} Replace \texttt{flask==3.0.2} with \texttt{flask} (unpinned), build a fresh environment, and note the version you get. Explain what could break six months from now, then pin it again.
\item
  \textbf{Extend.} Set up \texttt{pip-tools}: move the direct dependencies into \texttt{requirements.in}, generate \texttt{requirements.txt} with \texttt{pip-compile}, and read the comments it writes.
\item
  \textbf{Explain.} What is a transitive dependency, and why must it be pinned too?
\end{enumerate}

\noindent{}Solutions and hints are in the companion repository, in \texttt{exercises/}\allowbreak{}\texttt{chapter-12.md}. Try each exercise before reading its solution: the point is the thinking, not the answer.

\unnumberedsection{false}{What's Next}\label{12-dependency-declaration:whats-next}

Dependencies are declared. The next step is handling the configuration values that the application reads at runtime~--- moving every environment-specific value out of the source code and into a properly structured, environment-controlled configuration system.

\noindent{}→ \hyperref[ch:13-configuration-externalization]{Chapter 13}~--- Configuration Externalization
